\section{Background}
\label{sec:background}

\subsection{The Drake equation and what exoplanet science settled}

The Drake equation organises the expected number of communicative civilizations
as a product,
\begin{equation}
  N = R_* f_p n_e f_l f_i f_c L ,
  \label{eq:drake}
\end{equation}
and is best read as a map of uncertainty rather than a calculator
\citep{drake1992}. Exoplanet surveys have substantially constrained its leading
astrophysical factors: planets are common and potentially habitable planets are
no longer speculative \citep{borucki2016,batalha2014,dressing2015,kopparapu2013,seager2013}.
\Established

The consequence is that uncertainty has migrated to the right of
Eq.~\eqref{eq:drake}: to the emergence of life and intelligence, and above all to
$L$, the interval over which a civilization is detectable. \citet{prantzos2020}
shows explicitly that the $L$ factor dominates the inference. Our framework acts
on precisely this term, by refusing to treat it as a scalar.

\subsection{Great Filter and existential risk}

The Great Filter frames the problem as a bottleneck: some step between
non-living matter and a long-lived expansive civilization must be extremely
improbable \citep{hanson1998filter}. Whether that step is behind or ahead of us
carries obvious stakes, and the literature on existential risk explores the
latter possibility in detail \citep{bostrom2002,bostrom2013,ord2020,haggstrom2016}.
\citet{cirkovic2018} surveys the philosophical terrain, and
\citet{verendel2017} give a Bayesian treatment.

The Great Filter is a statement about probability mass, not about dynamics. It
does not, by itself, say how a civilization moves through the bottleneck, nor
what it looks like from outside while doing so. That is the gap a dynamical
model can address.

\subsection{Technosignatures and the search}

What could actually be detected has been studied since \citet{cocconi1959} and
\citet{dyson1960}. Energy-use signatures scale with the Kardashev classification
\citep{kardashev1964}; infrared searches for large-scale energy capture have
produced null results at survey scale \citep{wright2014ghat,griffith2015,garrett2015};
and the field has increasingly formalised how to reason about coverage of the
search space \citep{tarter2001,wright2018state,sheikh2020,wright2022haystack}.

Two points matter here. First, different signature classes correspond to
different civilizational behaviours --- waste heat, narrowband beacons,
propulsion, megastructures, atmospheric markers --- so ``detectability'' is not
one quantity. Second, the searches conducted so far cover a small fraction of
the parameter space, which is why $\Psearch$ must appear explicitly in any
honest inference \citep{wright2022haystack}. \Established

\subsection{Expansion, percolation and timing}

If interstellar settlement is feasible, the galaxy should be crossed quickly on
astronomical timescales \citep{hart1975,tipler1980,armstrong2013}, which sharpens
the paradox considerably. Percolation and agent-based treatments soften it by
allowing settlement to stall \citep{landis1998,forgan2009}. Anthropic and
hard-steps arguments address why we find ourselves early
\citep{carter1983,watson2008,kipping2020}, and the grabby-aliens analysis of
\citet{hanson2021grabby} makes the sharp observation that if loud civilizations
explain human earliness, then quiet ones must also be rare.

That last result is the most direct challenge to any low-observability account
of the Great Silence, including ours, and we return to it in
\S\ref{sec:discussion}.

\subsection{Recursive intelligence}

The possibility that capability improves the machinery of capability
has been discussed since \citet{good1965}, and elaborated with attention to
control and alignment by \citet{bostrom2014}, \citet{russell2019},
\citet{omohundro2008} and \citet{yudkowsky2008}. Quantum computation is relevant
to this framework only through this channel --- as one possible accelerant of
the recursive improvement loop \citep{shor1994,preskill2018,nielsen2010} --- and
we are explicit that its magnitude and even its sign of effect are unconstrained.
\Speculative

\subsection{The gap}

Across these literatures, detectability is treated as a static property: a
coefficient, a lifetime, or a survey limit. Where it is allowed to change, the
change is typically imposed as a chosen profile. What is missing is a
formulation in which observability has its own equation of motion, is coupled to
the same processes that drive capability and regulation, and is constrained to
remain physical. Supplying that, and testing it seriously, is the purpose of
this paper.
